\documentclass[11pt]{article}

\usepackage[a4paper,margin=1in]{geometry}
\usepackage{amsmath,amssymb,amsthm}
\usepackage[T1]{fontenc}
\usepackage{lmodern}
\usepackage[hidelinks]{hyperref}
\setlength{\emergencystretch}{3em}

\theoremstyle{plain}
\newtheorem{theorem}{Theorem}
\newtheorem{lemma}[theorem]{Lemma}
\newtheorem{proposition}[theorem]{Proposition}
\newtheorem{corollary}[theorem]{Corollary}
\theoremstyle{definition}
\newtheorem{definition}[theorem]{Definition}
\theoremstyle{remark}
\newtheorem{remark}[theorem]{Remark}

\newcommand{\conjid}{00000000088}
\newcommand{\Fu}{\mathcal{F}^{u}}
\newcommand{\Fs}{\mathcal{F}^{s}}
\newcommand{\HH}{H^{1}}
\newcommand{\ZZ}{\mathbb{Z}}
\newcommand{\QQ}{\mathbb{Q}}
\newcommand{\RR}{\mathbb{R}}
\newcommand{\CC}{\mathbb{C}}

\title{No pseudo-Anosov stretch factor is rational\\[0.4em]
\large The Last Math Competition, Conjecture \conjid}
\author{%
  Feiyu\thanks{Independent researcher.}
  \and
  Claude (Opus 5)\thanks{Anthropic.}%
}
\date{2026-09-12}

\begin{document}
\maketitle

\begin{abstract}
Conjecture \conjid{} asserts that every prime $p\ge2$ is the stretch factor of a
pseudo-Anosov element of some mapping class group. It is false --- and false for
every prime at once, because of a statement that is stronger and cleaner than the
negation: \emph{no pseudo-Anosov stretch factor is rational.} A stretch factor
$\lambda$ is an eigenvalue of the action of a homeomorphism on the first
cohomology of the double cover that orients the invariant foliations; that action
is an automorphism of a lattice, so its characteristic polynomial is monic with
constant term $\pm1$ and $\lambda$ is an algebraic unit. A rational algebraic unit
is $\pm1$, whereas $\lambda>1$. We give the argument in full, and we explain why
the Perron--Frobenius candidate matrix recorded in the research notes for this
conjecture cannot be repaired: the obstruction is not the determinant of a
transition matrix, and no choice of matrix, train track or surface avoids it. The
arithmetic half of the proof is formalized in Lean~4 against Mathlib; the
geometric half, which needs measured foliations, is classical and is given here by
hand.
\end{abstract}

\section{The conjecture as stated}

The text of conjecture \conjid{} reads:

\begin{quote}
\textbf{English.} Conjecture: For every prime $p \ge 2$ there exists a
pseudo-Anosov element in some mapping class group whose stretch factor is exactly
the integer $p$.
\end{quote}

\noindent
(The file \texttt{conjectures/\conjid.md} gives the same statement in Chinese; we
quote the English text only, so that this document compiles without a CJK font.)

\section{Reading of the statement}

The statement is unambiguous once three conventions are fixed. We fix them in the
way that is most generous to the conjecture: every reading we could give it is
refuted below.

\subsection{Surfaces and mapping classes}

Let $S$ be a surface of finite type: a closed surface with finitely many points
removed (punctures) or marked, possibly with finitely many boundary components.
$\mathrm{MCG}(S)$ is the group of isotopy classes of homeomorphisms of $S$
preserving the set of punctures and the boundary setwise. ``Some mapping class
group'' means $\mathrm{MCG}(S)$ for some such $S$; we allow $S$ to be orientable
or non-orientable, closed or not, and we place no bound on its genus or on the
number of punctures. Infinite-type surfaces are excluded, as they must be: the
Nielsen--Thurston classification, and with it the notion of a stretch factor, is a
finite-type statement.

\subsection{Pseudo-Anosov maps and stretch factors}

\begin{definition}\label{def:pa}
A homeomorphism $\varphi\colon S\to S$ is \emph{pseudo-Anosov} if there are two
transverse measured foliations $(\Fu,\mu^{u})$, $(\Fs,\mu^{s})$ on $S$ and a real
number $\lambda>1$ with
\[
  \varphi\cdot(\Fu,\mu^{u})=(\Fu,\lambda\mu^{u}),
  \qquad
  \varphi\cdot(\Fs,\mu^{s})=(\Fs,\lambda^{-1}\mu^{s}).
\]
The number $\lambda=\lambda(\varphi)$ is the \emph{stretch factor} (also
\emph{dilatation}) of $\varphi$. A mapping class is pseudo-Anosov if it has a
pseudo-Anosov representative; the stretch factor depends only on the class
\cite{FLP,FM}.
\end{definition}

Singularities are allowed: $\Fu$ and $\Fs$ have finitely many singular points with
$k\ge3$ prongs, and $k\ge1$ prongs at punctures and boundary components. This is
the definition in \cite{FLP,FM,Thurston88}, and the one under which the
Nielsen--Thurston classification is stated.

\subsection{Which number is ``the stretch factor''}\label{sec:convention}

Two normalizations appear in the literature: the stretch factor $\lambda$ of
Definition~\ref{def:pa}, which is $e^{\ell}$ for $\ell$ the translation length of
$\varphi$ on Teichmüller space; and the quasiconformal dilatation $K=\lambda^{2}$
of the Teichmüller representative. The conjecture is refuted under both, and more
generally under any convention that names $\lambda^{k}$ for a fixed integer
$k\ge1$: see Remark~\ref{rem:convention}. We work with $\lambda$.

\section{Result}

\begin{theorem}[Stretch factors are algebraic units]\label{thm:unit}
Let $\varphi$ be a pseudo-Anosov homeomorphism of a finite-type surface,
orientable or not. Then $\lambda(\varphi)$ is an algebraic unit: it is an
algebraic integer, and so is $\lambda(\varphi)^{-1}$.
\end{theorem}

This is classical --- it is the ``unit'' half of Fried's theorem that stretch
factors are bi-Perron algebraic units \cite{Fried}, and it is used as standard
equipment in the literature (e.g.\ \cite[\S5]{LeiningerReid}). We give a complete
proof in \S\ref{sec:proof} because it is what the refutation rests on.

\begin{theorem}[No rational stretch factors]\label{thm:irrational}
No pseudo-Anosov stretch factor is a rational number. In particular no integer
$n\ge2$ is a stretch factor.
\end{theorem}

\begin{corollary}\label{cor:main}
Conjecture \conjid{} is false, and it fails for \emph{every} prime $p\ge2$: for no
prime $p$ does any mapping class group contain a pseudo-Anosov element of stretch
factor $p$.
\end{corollary}

So the conjecture is not false by an exceptional prime or in a boundary case: the
class of numbers it asks for is disjoint from the class of numbers that can occur.
The arithmetic content of Theorems~\ref{thm:unit} and~\ref{thm:irrational} --- an
eigenvalue of a unimodular integer matrix that is rational equals $\pm1$ --- is
formalized in Lean~4 against Mathlib in \texttt{lean/Submission/Basic.lean}; see
\S\ref{sec:formal} for exactly what is and is not machine-checked.

\section{Proof}\label{sec:proof}

\subsection{The classical input}\label{sec:input}

We use four standard facts. Throughout, $\Sigma$ denotes the closed surface
underlying $S$, obtained by collapsing each boundary component to a puncture and
then filling in all punctures; a homeomorphism of $S$ induces a homeomorphism of
$\Sigma$ permuting the filled points.

\begin{itemize}
\item[(F1)] \textbf{(Orienting double cover.)} Let $\varphi$ be pseudo-Anosov on
$S$ with invariant foliations $\Fu,\Fs$. There is a double cover
$\pi\colon\widehat S\to S$, branched exactly over the singularities with an odd
number of prongs, such that $\pi^{*}\Fu$ and $\pi^{*}\Fs$ are orientable, and
$\varphi$ lifts to a homeomorphism $\widehat\varphi$ of $\widehat S$. The lift is
pseudo-Anosov with invariant foliations $\pi^{*}\Fu,\pi^{*}\Fs$ and
\[
  \lambda(\widehat\varphi)=\lambda(\varphi).
\]
(The cover is the one associated with the orientation character
$w\colon\pi_{1}(S\smallsetminus\mathrm{Sing})\to\ZZ/2$ of the tangent line field
of $\Fu$. Since $\varphi$ preserves $\Fu$, it preserves $w$, hence preserves
$\ker w$ --- a subgroup of index $2$, so normal, so unaffected by the basepoint
ambiguity of $\varphi_{*}$ --- and therefore lifts. The transverse measures pull
back, and $\pi\circ\widehat\varphi=\varphi\circ\pi$ gives
$\widehat\varphi\cdot\pi^{*}\mu^{u}=\pi^{*}(\varphi\cdot\mu^{u})
=\lambda\,\pi^{*}\mu^{u}$.) In the flat-geometry language this is the abelian
double cover on which the invariant quadratic differential becomes the square of
an abelian differential; see \cite[\S2]{Zorich}, \cite[Exp.~5, 9]{FLP},
\cite{Strenner}.

\item[(F2)] \textbf{(Oriented foliations are closed $1$-forms.)} An orientable
measured foliation on a closed surface $\Sigma$, together with a choice of leafwise
orientation, is the kernel foliation of a closed $1$-form $\omega$ with
$|\omega|=\mu$, vanishing exactly at the singularities; it therefore has a
well-defined class $[\mu]:=[\omega]\in\HH(\Sigma;\RR)$, and for every
homeomorphism $\psi$ of $\Sigma$,
\[
  [\psi\cdot\mu]=(\psi^{-1})^{*}[\mu].
\]
Note that orientability forces every singularity to have an even number of prongs,
so the foliation and the form extend across filled punctures and the passage from
$S$ to $\Sigma$ is harmless. See \cite[\S2]{Zorich}, \cite{FLP}.

\item[(F3)] \textbf{(Nondegeneracy.)} If $\Fu,\Fs$ is the transverse pair of a
pseudo-Anosov map and both are orientable, then the $1$-forms of (F2) satisfy
\[
  \int_{\Sigma}\omega^{u}\wedge\omega^{s}=\mathrm{Area}>0 ,
\]
because $\omega^{u}+i\,\omega^{s}$ is an abelian differential whose horizontal and
vertical foliations are $\Fu$ and $\Fs$ and $\omega^{u}\wedge\omega^{s}$ is its
area form. Hence $[\mu^{u}]\smile[\mu^{s}]\ne0$ in $H^{2}(\Sigma;\RR)$, and in
particular
\[
  [\mu^{u}]\ne0 .
\]

\item[(F4)] \textbf{(Homological actions are unimodular.)} For a closed surface
$\Sigma$, $\HH(\Sigma;\ZZ)\cong\ZZ^{m}$ with $m=2g$, and a homeomorphism $\psi$
acts on it by a group automorphism. In a basis, $\psi^{*}$ is a matrix
$M\in\mathrm{GL}_{m}(\ZZ)$, so $\det M=\pm1$.
\end{itemize}

\subsection{Orientable foliations}

\begin{proposition}\label{prop:orientable}
Let $\varphi$ be a pseudo-Anosov homeomorphism of a finite-type orientable surface
$S$ whose invariant foliations are orientable, and let $\Sigma$ be the closed
surface obtained by filling the punctures. Then there are $M\in
\mathrm{GL}_{m}(\ZZ)$ and $\varepsilon\in\{\pm1\}$ such that $\varepsilon\lambda$
is an eigenvalue of $M$, where $\lambda=\lambda(\varphi)$.
\end{proposition}

\begin{proof}
Fix a leafwise orientation of $\Fu$ and let $\omega^{u}$ be the closed $1$-form of
(F2), so that $[\mu^{u}]=[\omega^{u}]\in\HH(\Sigma;\RR)$. The homeomorphism
$\varphi$ either preserves or reverses the chosen orientation; write
$\varepsilon=+1$ in the first case and $\varepsilon=-1$ in the second. Applying
(F2) to $\psi=\varphi$ and using $\varphi\cdot\mu^{u}=\lambda\mu^{u}$,
\[
  (\varphi^{-1})^{*}[\mu^{u}]
  =[\varphi\cdot\mu^{u}]
  =\varepsilon\lambda\,[\mu^{u}] .
\]
By (F3), $[\mu^{u}]\ne0$, so $\varepsilon\lambda$ is an eigenvalue of the
$\RR$-linear map $(\varphi^{-1})^{*}$ on $\HH(\Sigma;\RR)$. That map preserves the
lattice $\HH(\Sigma;\ZZ)$ and is invertible on it, so by (F4) its matrix $M$ in an
integral basis lies in $\mathrm{GL}_{m}(\ZZ)$, and
$\HH(\Sigma;\RR)=\HH(\Sigma;\ZZ)\otimes\RR$ identifies the two actions.
\end{proof}

\subsection{The general case}\label{sec:general}

\begin{proposition}\label{prop:general}
Let $\varphi$ be a pseudo-Anosov homeomorphism of any finite-type surface,
orientable or not, and $\lambda=\lambda(\varphi)$. Then there are
$M\in\mathrm{GL}_{m}(\ZZ)$ and $\varepsilon\in\{\pm1\}$ such that
$\varepsilon\lambda$ is an eigenvalue of $M$.
\end{proposition}

\begin{proof}
First let $S$ be orientable. By (F1) pass to the orienting double cover: we obtain
a pseudo-Anosov $\widehat\varphi$ on the finite-type orientable surface $\widehat
S$ with orientable invariant foliations and
$\lambda(\widehat\varphi)=\lambda(\varphi)=\lambda$. Proposition
\ref{prop:orientable} applied to $\widehat\varphi$ gives the claim.

If $S$ is non-orientable, let $\rho\colon\widetilde S\to S$ be its orientation
double cover. It is canonical, so every homeomorphism of $S$ lifts; the lift
$\widetilde\varphi$ of $\varphi$ preserves the pulled-back measured foliations and
scales their measures by the same factors, so $\widetilde\varphi$ is
pseudo-Anosov on the orientable surface $\widetilde S$ with
$\lambda(\widetilde\varphi)=\lambda$ (see \cite{Strenner}). The orientable case
just proved applies to $\widetilde\varphi$.
\end{proof}

\begin{remark}
The torus needs no special treatment. If one admits the Anosov elements of
$\mathrm{MCG}(T^{2})\cong\mathrm{SL}_{2}(\ZZ)$ as ``pseudo-Anosov'', their stretch
factors are by definition eigenvalues of matrices in $\mathrm{SL}_{2}(\ZZ)$, so the
conclusion of Proposition~\ref{prop:general} holds for them verbatim, and
everything below applies.
\end{remark}

\subsection{Eigenvalues of unimodular integer matrices}

\begin{lemma}\label{lem:unit}
Let $M\in\mathrm{GL}_{m}(\ZZ)$ and let $\alpha\in\CC$ be an eigenvalue of $M$.
Then $\alpha$ is an algebraic unit.
\end{lemma}

\begin{proof}
Let $\chi(X)=X^{m}+c_{m-1}X^{m-1}+\dots+c_{1}X+c_{0}\in\ZZ[X]$ be the
characteristic polynomial of $M$; it is monic, and
$c_{0}=\chi(0)=\det(-M)=(-1)^{m}\det M=\pm1$. Thus $\alpha$ is an algebraic
integer. Moreover $\chi(\alpha)=0$ gives, after dividing by $\alpha\ne0$,
\[
  \alpha^{-1}
  =-c_{0}^{-1}\bigl(\alpha^{m-1}+c_{m-1}\alpha^{m-2}+\dots+c_{1}\bigr)
  =\mp\bigl(\alpha^{m-1}+c_{m-1}\alpha^{m-2}+\dots+c_{1}\bigr),
\]
a $\ZZ$-polynomial expression in the algebraic integer $\alpha$; hence
$\alpha^{-1}$ is an algebraic integer as well.
\end{proof}

\begin{proof}[Proof of Theorem~\ref{thm:unit}]
By Proposition~\ref{prop:general}, $\varepsilon\lambda$ is an eigenvalue of some
$M\in\mathrm{GL}_{m}(\ZZ)$, so $\varepsilon\lambda$ is an algebraic unit by
Lemma~\ref{lem:unit}. Algebraic units are closed under negation, so $\lambda$ is
one.
\end{proof}

\subsection{Arithmetic}

\begin{lemma}\label{lem:rational}
A rational algebraic unit is $1$ or $-1$.
\end{lemma}

\begin{proof}
Let $\alpha\in\QQ$ be an algebraic unit. An algebraic integer in $\QQ$ is in $\ZZ$
($\ZZ$ is integrally closed in $\QQ$), so $\alpha\in\ZZ$ and $\alpha^{-1}\in\ZZ$.
Then $\alpha\cdot\alpha^{-1}=1$ with both factors integers, so $\alpha$ divides
$1$ in $\ZZ$, i.e.\ $\alpha=\pm1$.
\end{proof}

Equivalently, and this is the form that is formalized: if $\alpha\in\ZZ$ is a root
of a monic $\chi\in\ZZ[X]$, then $\alpha\mid\chi(\alpha)-\chi(0)$, so
$\alpha\mid\chi(0)=\pm1$.

\begin{proof}[Proof of Theorem~\ref{thm:irrational}]
Suppose $\lambda=\lambda(\varphi)\in\QQ$ for some pseudo-Anosov $\varphi$. By
Theorem~\ref{thm:unit}, $\lambda$ is an algebraic unit, so $\lambda=\pm1$ by
Lemma~\ref{lem:rational}. But $\lambda>1$ by Definition~\ref{def:pa}. This is a
contradiction, so no stretch factor is rational; integers $n\ge2$ are in
particular rational.
\end{proof}

\begin{proof}[Proof of Corollary~\ref{cor:main}]
A prime $p\ge2$ is an integer $\ge2$, hence is not a stretch factor by
Theorem~\ref{thm:irrational}. Since $p$ was an arbitrary prime, the conjecture
fails for every one of them.
\end{proof}

\begin{remark}[the refutation does not depend on the convention]
\label{rem:convention}
Suppose ``stretch factor'' is read as $\lambda^{k}$ for some fixed integer $k\ge1$
--- for instance $k=2$, the quasiconformal dilatation $K=\lambda^{2}$ of the
Teichmüller representative. If $\lambda$ is an algebraic unit then so is
$\lambda^{k}$, and $\lambda^{k}>1$; so by Lemma~\ref{lem:rational} the number
$\lambda^{k}$ is again never rational. Every reading of the conjecture of this
shape is false. Likewise, asking for $\lambda=\sqrt p$ (so that $K=p$) fails: the
minimal polynomial of $\sqrt p$ is $X^{2}-p$, whose constant term is not $\pm1$, so
$\sqrt{p}$ is not a unit.
\end{remark}

\begin{remark}[what the conjecture asks for, in one line]
The stretch factors of pseudo-Anosov maps are contained in the bi-Perron algebraic
units \cite{Fried}. A prime $p$ is bi-Perron --- trivially, having no other
conjugates --- but it is not a unit. That single missing condition is what the
conjecture overlooks, and it is fatal.
\end{remark}

\section{Why the Perron--Frobenius candidate cannot be repaired}
\label{sec:candidate}

The research notes accompanying this conjecture propose, for each prime $p\ge3$,
the positive integer matrix
\[
  M_{p}=\begin{pmatrix}2&1\\p-2&p-1\end{pmatrix},
  \qquad
  \chi_{M_{p}}(X)=X^{2}-(p+1)X+p=(X-p)(X-1),
\]
which is primitive with Perron--Frobenius eigenvalue exactly $p$, and observe that
the associated two-edge substitution $e_{1}\mapsto e_{1}^{2}e_{2}^{p-2}$,
$e_{2}\mapsto e_{1}e_{2}^{p-1}$ induces $M_{p}$ on the first homology of the
rank-two graph, where $\det M_{p}=p\ne\pm1$ rules out a homotopy equivalence. Both
computations are correct; they are re-checked formally in the Lean development
(\texttt{cand\_charpoly}, \texttt{cand\_det}). Three comments place them in
context.

\smallskip
\noindent\textbf{1. The Perron--Frobenius side has no obstruction at all.} By
Lind's theorem \cite{Lind}, the spectral radii of primitive non-negative integer
matrices are exactly the Perron numbers, and every integer $p\ge2$ is a Perron
number. Matrices with Perron root $p$ therefore exist in every sufficiently large
size and in great abundance; $M_{p}$ is one of infinitely many. A search over transition matrices therefore never
fails for arithmetic reasons --- and, by Theorem~\ref{thm:unit}, never succeeds
for geometric ones.

\smallskip
\noindent\textbf{2. The notes' determinant obstruction does not generalize by
itself.} A transition matrix of an invariant train track \cite{BH} counts, entry by
entry, how many times the image of one branch runs over another ---
\emph{unsigned} crossing numbers.
The homological action counts them with signs. The two matrices agree only when no
cancellation occurs, which is what happens for the positive substitution above and
is why the argument in the notes is valid there. In general the transition matrix
is not the homological matrix, and unimodularity is a constraint on the latter, not
the former. So the obstruction found in the notes is genuine for that candidate but
is not the general mechanism, and one cannot conclude from it that higher-rank
train tracks are barred.

\smallskip
\noindent\textbf{3. The general mechanism is Theorem~\ref{thm:unit}, and it closes
the problem.} Whatever the transition matrix, a pseudo-Anosov map has, after the
orienting double cover, a \emph{homological} action that is unimodular and has
$\pm\lambda$ among its eigenvalues. So the conclusion of the notes --- ``no such
realization has been established here'' --- can be strengthened to: no such
realization exists, for any train track of any size on any surface, for any prime
$p$. The correct reading of the notes' determinant computation is that it is the
shadow, in the simplest case, of an obstruction that lives one level up.

\section{Formalization}\label{sec:formal}

The accompanying Lean~4 project (\texttt{lean/}, namespace
\texttt{Conjecture00000000088}, Mathlib \texttt{v4.33.1}) formalizes the
arithmetic half of the proof --- everything from
Proposition~\ref{prop:general} onwards. Measured foliations, train tracks and
Teichmüller theory are not in Mathlib, so \S\ref{sec:input}--\S\ref{sec:general}
(the facts (F1)--(F4) and Propositions~\ref{prop:orientable},
\ref{prop:general}) are \emph{not} machine-checked; they are classical, and are
proved above from standard references. We state this division plainly rather than
formalizing a restatement of the conjecture that would hide it.

The interface between the two halves is the predicate
\begin{quote}
\texttt{IsUnimodularEigenvalue lam} $:=$ there are $m$, an integer matrix
$A\in\ZZ^{m\times m}$ with \texttt{IsUnit A.det}, and a nonzero real vector $v$
with $A v=\texttt{lam}\cdot v$,
\end{quote}
which is exactly the conclusion of Proposition~\ref{prop:general} (up to the sign
$\varepsilon$, which is immaterial: $M$ may be replaced by $-M$, still
unimodular). What is proved in Lean:

\begin{itemize}
  \item \texttt{int\_root\_dvd\_coeff\_zero} --- an integer root of an integer
    polynomial divides its constant coefficient.
  \item \texttt{isUnit\_charpoly\_coeff\_zero} --- for a unimodular integer
    matrix, $\chi(0)=\pm\det$ is a unit (Mathlib's
    \texttt{Matrix.det\_eq\_sign\_charpoly\_coeff}).
  \item \texttt{eq\_one\_or\_neg\_one\_of\_charpoly\_eval} --- an integer
    eigenvalue of a unimodular integer matrix is $\pm1$; this is
    Lemma~\ref{lem:unit} plus Lemma~\ref{lem:rational} in the integral case.
  \item \texttt{exists\_charpoly\_map\_eval\_eq\_zero} --- a real eigenvector for
    \texttt{lam} makes the characteristic polynomial vanish at \texttt{lam}
    (Mathlib's \texttt{Matrix.exists\_mulVec\_eq\_zero\_iff} and
    \texttt{Matrix.eval\_charpoly}).
  \item \texttt{rat\_eq\_one\_or\_neg\_one} --- a \emph{rational} number
    satisfying \texttt{IsUnimodularEigenvalue} is $1$ or $-1$. The descent from
    $\RR$ to $\QQ$ to $\ZZ$ uses that $\ZZ$ is integrally closed in $\QQ$
    (\texttt{IsIntegrallyClosed.isIntegral\_iff}), the characteristic polynomial
    being monic.
  \item \texttt{irrational\_of\_isUnimodularEigenvalue} --- hence any such number
    $>1$ is irrational: Theorem~\ref{thm:irrational}, modulo the geometry.
  \item \texttt{prime\_not\_isUnimodularEigenvalue} --- hence no prime satisfies
    the predicate: Corollary~\ref{cor:main}, modulo the geometry.
  \item \texttt{goldenRatio\_isUnimodularEigenvalue} --- the golden ratio does
    satisfy it, via $\bigl(\begin{smallmatrix}1&1\\1&0\end{smallmatrix}\bigr)$. The
    theorems above are therefore not vacuous: what they exclude is rationality and
    nothing more.
  \item \texttt{cand\_charpoly}, \texttt{cand\_det},
    \texttt{cand\_not\_unimodular} --- the two computations of the research notes
    for $M_{p}$, and the conclusion that $M_{p}$ is not unimodular for $p\ge2$.
  \item \texttt{repaired\_det}, \texttt{repaired\_charpoly},
    \texttt{repaired\_isUnimodularEigenvalue} --- the repaired candidate $M_{p}'$
    of \S\ref{sec:repair}: it is unimodular, its characteristic polynomial is
    $X^{2}-pX+1$, and its larger root passes, for every integer $p\ge2$, the test
    that $p$ itself fails.
\end{itemize}

Primality is Mathlib's \texttt{Nat.Prime} and irrationality is Mathlib's
\texttt{Irrational}; no standard notion is redefined. The project builds with
\texttt{lake build} on \texttt{leanprover/lean4:v4.33.1}. It contains no
\texttt{sorry}, no \texttt{admit} and no \texttt{native\_decide}, and introduces no
axioms of its own: \texttt{\#print axioms} reports exactly \texttt{propext},
\texttt{Classical.choice} and \texttt{Quot.sound} for each theorem above.

\section{What is true instead}

\subsection{The nearest true statement}\label{sec:repair}

The obstruction is exactly unimodularity, so the natural repair keeps the integer
$p$ but moves it from the eigenvalue to the trace. Consider
\[
  M_{p}'=\begin{pmatrix}1&1\\p-2&p-1\end{pmatrix}
  \qquad\text{(the notes' } M_{p}\text{ with its top-left entry lowered by }1),
\]
so that $\det M_{p}'=(p-1)-(p-2)=1$ and $\operatorname{tr}M_{p}'=p$, whence
$\chi_{M_{p}'}(X)=X^{2}-pX+1$ and
\[
  \lambda_{p}=\frac{p+\sqrt{p^{2}-4}}{2},\qquad \lambda_{p}+\lambda_{p}^{-1}=p .
\]
For $p\ge3$ this matrix is hyperbolic, and under the isomorphism
$\mathrm{MCG}(S_{1,1})\cong\mathrm{SL}_{2}(\ZZ)$ it is a pseudo-Anosov mapping
class of the once-punctured torus with stretch factor $\lambda_{p}$
\cite[\S13.1]{FM}; equivalently it is an Anosov diffeomorphism of $T^{2}$. Thus:

\begin{quote}
For every prime $p\ge3$ there is a pseudo-Anosov mapping class whose stretch
factor $\lambda$ satisfies $\lambda+\lambda^{-1}=p$.
\end{quote}

This is true, elementary, and is what the conjecture should have asked. It also
shows the price of unimodularity precisely: for $p=3$ one gets
$\lambda=(3+\sqrt5)/2\approx2.618$ rather than $3$. The prime $p=2$ is genuinely
exceptional even for the repaired statement --- $\operatorname{tr}=2$ is parabolic,
$\lambda=1$, the mapping class is a Dehn twist and is reducible, not pseudo-Anosov.

\subsection{Where the boundary of the possible lies}

Fried \cite{Fried} proved that stretch factors are bi-Perron algebraic units and
conjectured the converse: that every bi-Perron algebraic unit has a power that is a
stretch factor. That conjecture remains open, and recent work suggests it is
false in a strong quantitative sense --- asymptotically almost all bi-Perron units
of degree at most $2n$ are not stretch factors on a closed genus-$n$ surface
\cite{BRW}. So the exact set of stretch factors is not known; what is known, and
all that is needed here, is that it contains no rational number.

\subsection{The contrast that makes the conjecture tempting}

Integer growth rates are perfectly common one dimension down. By Lind's theorem
\cite{Lind} every Perron number, $p$ included, is the spectral radius of a
primitive non-negative integer matrix; and by Thurston's classification of
one-dimensional entropies \cite{ThurstonEntropy}, $\log p$ is the topological
entropy of a postcritically finite interval map. What fails for surfaces is not the
Perron--Frobenius arithmetic but the extra rigidity a surface imposes: the
transverse measure of an invariant foliation is a cohomology class of the surface,
and a homeomorphism moves cohomology by a unimodular integer matrix.

\begin{thebibliography}{99}

\bibitem{BRW} H.~Baik, A.~Rafiqi, C.~Wu,
\emph{Is a typical bi-Perron algebraic unit a pseudo-Anosov dilatation?},
arXiv:1612.03084.

\bibitem{BH} M.~Bestvina, M.~Handel,
\emph{Train-tracks for surface homeomorphisms},
Topology \textbf{34} (1995), 109--140.

\bibitem{FM} B.~Farb, D.~Margalit,
\emph{A Primer on Mapping Class Groups},
Princeton Mathematical Series \textbf{49}, Princeton University Press, 2012.

\bibitem{FLP} A.~Fathi, F.~Laudenbach, V.~Poénaru,
\emph{Travaux de Thurston sur les surfaces},
Astérisque \textbf{66--67} (1979);
English translation: \emph{Thurston's Work on Surfaces},
Mathematical Notes \textbf{48}, Princeton University Press, 2012.

\bibitem{Fried} D.~Fried,
\emph{Growth rate of surface homeomorphisms and flow equivalence},
Ergodic Theory Dynam. Systems \textbf{5} (1985), 539--563.

\bibitem{LeiningerReid} C.~J.~Leininger, A.~W.~Reid,
\emph{Pseudo-Anosov homeomorphisms not arising from branched covers},
Groups Geom. Dyn. \textbf{14} (2020), 151--175.

\bibitem{Lind} D.~A.~Lind,
\emph{The entropies of topological Markov shifts and a related class of algebraic
integers},
Ergodic Theory Dynam. Systems \textbf{4} (1984), 283--300.

\bibitem{Strenner} B.~Strenner,
\emph{Lifts of pseudo-Anosov homeomorphisms of nonorientable surfaces have
vanishing SAF invariant},
Math. Res. Lett. \textbf{25} (2018), 677--685.

\bibitem{Thurston88} W.~P.~Thurston,
\emph{On the geometry and dynamics of diffeomorphisms of surfaces},
Bull. Amer. Math. Soc. (N.S.) \textbf{19} (1988), 417--431.

\bibitem{ThurstonEntropy} W.~P.~Thurston,
\emph{Entropy in dimension one}, in:
Frontiers in Complex Dynamics, Princeton University Press, 2014, 339--384.

\bibitem{Zorich} A.~Zorich,
\emph{Flat surfaces}, in:
Frontiers in Number Theory, Physics and Geometry I, Springer, 2006, 437--583.

\end{thebibliography}

\end{document}
